% =============================================================================
% Volume II: Course Overview — Day 1 Slide Deck
% =============================================================================
% This deck introduces the advanced course: Machine Learning Systems at Scale.
% Presented on the FIRST DAY to set context, scope, and expectations.
%
% IMAGES:
%   - single-to-fleet.pdf    - c3-taxonomy.pdf
%   - fleet-stack.pdf         - course-roadmap.pdf
%   - scale-numbers.pdf
% =============================================================================
\documentclass[aspectratio=169, 12pt]{beamer}
\usepackage{../../assets/beamerthememlsys}

\mlsyssetup{
  volume       = {Volume II},
  chapter      = {Course Overview},
  logo         = {../../assets/img/logo-mlsysbook.png},
  instlogo     = {../../assets/img/logo-harvard.png},
  chaptertitle = {ML Systems at Scale},
}

% --- Fonts ---
\usepackage{fontspec}
\setsansfont{Helvetica Neue}[
  BoldFont={Helvetica Neue Bold},
  ItalicFont={Helvetica Neue Italic},
  BoldItalicFont={Helvetica Neue Bold Italic},
]
\setmonofont{JetBrains Mono}[Scale=0.85]

% --- Packages ---
\usepackage{booktabs}
\usepackage{amsmath}

% --- Image paths ---
\graphicspath{
  {images/}
}

% --- Chapter-specific macros ---
\newcommand{\Cthree}{C\textsuperscript{3}}

% --- Helper: safe image include ---
\newcommand{\safeimg}[2][width=\textwidth,keepaspectratio]{%
  \IfFileExists{images/#2}{\includegraphics[#1]{#2}}{%
    \IfFileExists{#2}{\includegraphics[#1]{#2}}{%
      \fbox{\parbox[c][2.5cm][c]{0.85\linewidth}{\centering\footnotesize\textcolor{midgray}{[Missing image]}}}%
    }%
  }%
}

% --- Section count for navigation (must match actual \section{} count) ---
\setcounter{mlsystotalsections}{8}

\title{ML Systems at Scale}
\author{Vijay Janapa Reddi}
\institute{Harvard University}
\date{}

\begin{document}

% =============================================================================
% TITLE SLIDE
% =============================================================================
\mlsystitle{ML Systems at Scale}{The Fleet Is the Computer}{placeholder.pdf}

% =============================================================================
% WELCOME & COURSE IDENTITY
% =============================================================================
\section{Welcome}

\begin{frame}{Welcome to Volume II}
\note{
% -- LINK: What prior concept connects to this slide
Volume I taught the single-machine mental model: one node, 1--8 GPUs,
shared memory, PCIe/NVLink. This slide establishes that everything
students learned still applies --- but a new scale axis changes the rules.

% -- NARRATE: What to SAY while showing this slide
Point to the comparison table on the right: ``Every row flips when you
cross the node boundary. The bus becomes a network. Shared memory becomes
message passing. Rare failures become daily certainty.'' Use the tagline:
``The fleet is the computer'' --- echoing Sun Microsystems' ``The network
is the computer,'' but applied to ML clusters.

% -- ENGAGE: Specific question, prediction, or task for THIS slide
Ask: ``How many of you have SSH'd into a multi-GPU cluster?''
Hands up. Then: ``How many have debugged a training job that stalled
because one of 1,000 GPUs went silent?'' The gap between those two
counts is what this course fills.

% -- WARN: What students will get wrong on THIS topic
Common error: students assume ``distributed = just add more GPUs.''
Correct framing: crossing a node boundary changes failure modes,
communication patterns, and programming models qualitatively.
IF STUCK: Ask them to compare restarting a crashed browser tab vs.\
restarting one node in a 1,000-node synchronized training job.

% -- FLEX: [CORE] or [OPTIONAL] + contingency
[CORE] This is the opening frame --- it sets the tone for the entire course.
IF AHEAD: Ask what specific Vol I concept they found most surprising;
connect it to the fleet version.
IF SHORT: Skip the hand-raise, go straight to the table walkthrough.
}

\small
\begin{columns}[T]
  \begin{column}{0.55\textwidth}
    \textbf{Machine Learning Systems at Scale}

    \vspace{0.2cm}
    Volume I taught you how \emph{one machine} runs ML.

    \vspace{0.1cm}
    Volume II teaches you how \emph{thousands of machines} coordinate to train and serve the world's largest models.

    \vspace{0.2cm}
    \begin{mlsyscard}{crimson}
    \textbf{The Fleet Is the Computer}\\[0.1cm]
    {\footnotesize No single GPU trains GPT-4. No single server serves
    a billion users. The \emph{fleet} --- thousands of accelerators connected
    by network fabric --- is the unit of design.}
    \end{mlsyscard}
  \end{column}
  \begin{column}{0.42\textwidth}
    \vspace{0.3cm}
    \renewcommand{\arraystretch}{1.2}
    {\scriptsize
    \begin{tabular}{@{}lll@{}}
      \toprule
      & \textbf{Vol I} & \textbf{Vol II} \\
      \midrule
      Scope   & 1 machine     & 1,000+ machines \\
      Bus     & PCIe/NVLink   & IB/Ethernet \\
      Memory  & Shared        & Distributed \\
      Failure & Rare event    & Daily certainty \\
      API     & DataParallel  & torch.distributed \\
      \bottomrule
    \end{tabular}
    }
  \end{column}
\end{columns}

\end{frame}

% =============================================================================
\section{Single Node to Fleet}
% =============================================================================

\begin{frame}{From Single Node to Fleet}
\note{
% -- LINK: What prior concept connects to this slide
The Welcome slide introduced the fleet concept verbally. This diagram
makes the transition visual --- left side is Vol I territory, right side
is Vol II.

% -- NARRATE: What to SAY while showing this slide
Walk left-to-right through the diagram: ``On the left, one node ---
everything connected by NVLink at 900 GB/s, failures are rare, latency
is nanoseconds. Cross the dotted line to the right: the bus becomes
InfiniBand at 50 GB/s, latency jumps to microseconds, and with 10,000
GPUs, one fails every few hours. The network replaces the memory bus as
the critical interconnect.''

% -- ENGAGE: Specific question, prediction, or task for THIS slide
Ask: ``What happens to your training job when one of 10,000 GPUs dies?''
Give 15 seconds of think time. Expected answer: the entire job stalls
or crashes --- which is exactly why fault tolerance is first-class.

% -- WARN: What students will get wrong on THIS topic
Students will underestimate the latency jump: nanoseconds to microseconds
sounds small, but it is a 1,000x increase. At 350 GB of gradients per
step, that 1,000x turns into minutes of synchronization overhead per hour.
IF STUCK: Compare it to a highway going from 300 mph to 0.3 mph at a
toll booth --- the booth is the node boundary.

% -- FLEX: [CORE] or [OPTIONAL] + contingency
[CORE] This diagram anchors the entire course narrative.
IF AHEAD: ``At what fleet size does the probability of zero failures
during a 1-hour training window drop below 50\%?''
IF SHORT: Point to the diagram, read the insight callout, move on.
}

% --- Layout: FULL-WIDTH diagram ---
\centering
\safeimg[width=0.92\textwidth,height=4.8cm,keepaspectratio]{single-to-fleet.pdf}

\vspace{0.1cm}
\mlsysinsight{Key shift:}{The network is the bus, and \textbf{failure is a statistical certainty}.}

\end{frame}

\begin{frame}{Why Scale Changes Everything}
\note{
% -- LINK: What prior concept connects to this slide
The previous diagram showed the physical transition. This slide names
the three qualitative changes that make fleet engineering a different
discipline from single-node optimization.

% -- NARRATE: What to SAY while showing this slide
Point to each card in order. ``First, communication dominates ---
at 10,000 GPUs, AllReduce takes longer than the forward pass. Second,
failure is routine --- 10,000 GPUs times 100,000h MTBF equals one failure
every 10 hours. Third, emergent behavior --- stragglers, hot spots, and
cascading failures that no single component predicts.''
ANALOGY: ``A single car can break down. A fleet of 10,000 taxis has
at least one broken down at any moment --- and every taxi must wait for
every other taxi before the next fare.''

% -- ENGAGE: Specific question, prediction, or task for THIS slide
Ask: ``Which of these three surprises you most?'' Cold-call one student.
Most will say failure frequency --- use that to preview the reliability
math coming later.

% -- WARN: What students will get wrong on THIS topic
Students will treat these as independent problems. In reality they
interact: a straggler (emergent behavior) that triggers a timeout
(failure) during an AllReduce (communication) cascades across all three.
IF STUCK: Walk through a concrete cascade: slow GPU triggers BSP
barrier timeout, which triggers checkpoint, which stalls all 10,000 GPUs.

% -- FLEX: [CORE] or [OPTIONAL] + contingency
[CORE] These three categories structure the entire course.
IF AHEAD: ``Can you think of a fourth qualitative change at scale?''
(Answer: cost --- 1\% inefficiency at 10,000 GPUs is millions of dollars.)
IF SHORT: Read the three card titles and the bottom-line callout, skip
the analogy.
}

\small
\begin{columns}[T]
  \begin{column}{0.30\textwidth}
    \begin{mlsyscard}{computestroke}
    \textbf{Communication Dominates}\\[0.1cm]
    {\footnotesize At 10,000 GPUs, AllReduce takes longer than the forward pass. The network is the bottleneck, not the accelerator.}
    \end{mlsyscard}
  \end{column}
  \begin{column}{0.30\textwidth}
    \begin{mlsyscard}{errorstroke}
    \textbf{Failure is Routine}\\[0.1cm]
    {\footnotesize 10,000 GPUs $\times$ 100,000h MTBF = one failure every 10 hours. Checkpointing and recovery are first-class concerns.}
    \end{mlsyscard}
  \end{column}
  \begin{column}{0.30\textwidth}
    \begin{mlsyscard}{routingstroke}
    \textbf{Behavior Emerges}\\[0.1cm]
    {\footnotesize Stragglers, hot spots, and cascading failures. No single component determines system behavior.}
    \end{mlsyscard}
  \end{column}
\end{columns}

\vspace{0.2cm}
\centering
\mlsysalert{Bottom line:}{Scale is not ``more of the same'' --- it is a \textbf{qualitative} change in engineering.}

\end{frame}

% =============================================================================
\section{The \Cthree{} Framework}
% =============================================================================

\begin{frame}{The \Cthree{} Taxonomy: Compute, Communication, Coordination}
\note{
% -- LINK: What prior concept connects to this slide
The previous slide named three qualitative changes at scale. The C3
taxonomy formalizes these into a diagnostic framework --- it replaces
the single-node D-A-M taxonomy with a fleet-scale lens.

% -- NARRATE: What to SAY while showing this slide
Point to the diagram: ``Every performance problem in a fleet traces to
one of these three axes. Compute: are the GPUs doing useful math?
Communication: how fast can gradients and activations move? Coordination:
who decides what runs where, when to checkpoint, how to handle stragglers?''
Draw the parallel explicitly: ``Vol I asked `is this workload Data-bound,
Algorithm-bound, or Memory-bound?' Vol II asks `is the fleet bottlenecked
on Compute, Communication, or Coordination?'''

% -- ENGAGE: Specific question, prediction, or task for THIS slide
Ask: ``If training stalls for 30 seconds every hour, which C is the
culprit?'' Give 15 seconds. Cold-call. Expected answer: Coordination ---
likely synchronous checkpointing or straggler mitigation, not raw
compute or bandwidth.

% -- WARN: What students will get wrong on THIS topic
Students will conflate Communication and Coordination. Communication is
moving bytes (AllReduce, gradient sync). Coordination is making decisions
(scheduling, checkpointing, barrier management). A straggler that slows
AllReduce is a Coordination problem manifesting through Communication.
IF STUCK: ``Communication is the pipe. Coordination is the traffic cop.''

% -- FLEX: [CORE] or [OPTIONAL] + contingency
[CORE] C3 is the diagnostic backbone of the entire volume.
IF AHEAD: ``Can a problem be bottlenecked on two C's simultaneously?
Give an example.'' (Yes: gradient compression reduces Communication but
adds Compute overhead for encode/decode.)
IF SHORT: Show diagram, ask the 30-second stall question, move on.
}

% --- Layout: FULL-WIDTH diagram ---
\centering
\safeimg[width=0.92\textwidth,height=4.8cm,keepaspectratio]{c3-taxonomy.pdf}

\vspace{0.1cm}
\mlsysconcept{Diagnostic:}{``Is the fleet bottlenecked on \textcolor{computestroke}{Compute}, \textcolor{datastroke}{Communication}, or \textcolor{routingstroke}{Coordination}?''}

\end{frame}

\begin{frame}{\Cthree{} in Practice}
\note{
% -- LINK: What prior concept connects to this slide
The previous slide defined C3 abstractly. This table maps real fleet
symptoms to C3 axes, showing the framework in diagnostic action.

% -- NARRATE: What to SAY while showing this slide
Walk through the table row by row. ``Low GPU utilization? That is
Compute --- memory-bound kernels, fix with operator fusion. Throughput
plateau? Communication --- AllReduce saturated, fix with gradient
compression. Periodic 30-second stalls? Coordination --- synchronous
checkpointing, fix with async checkpointing.'' Emphasize the pattern:
diagnosis precedes optimization.

% -- ENGAGE: Specific question, prediction, or task for THIS slide
Cover the ``Response'' column with your hand. Point to ``Throughput
variance'' and ask: ``Which C axis and what would you do?'' Give 20
seconds. Expected: Coordination (straggler nodes), response is straggler
mitigation or redundant computation.

% -- WARN: What students will get wrong on THIS topic
Students will jump to solutions before diagnosing. ``Just buy faster
GPUs'' is the default instinct. The table shows that 4 of 6 symptoms
are NOT solved by faster hardware --- they require algorithmic or
systems-level interventions.
IF STUCK: Ask which rows would NOT be helped by doubling GPU TFLOPS.

% -- FLEX: [OPTIONAL] + contingency
[OPTIONAL] The C3 concept was introduced on the previous slide.
IF SHORT: Show the table briefly, read the insight callout, move on.
IF AHEAD: Ask students to propose a 7th row with a novel symptom.
}

\scriptsize
\renewcommand{\arraystretch}{1.15}
\begin{tabular}{@{}llll@{}}
  \toprule
  \textbf{Symptom} & \textbf{\Cthree{} Axis} & \textbf{Root Cause} & \textbf{Response} \\
  \midrule
  Low GPU utilization       & \textcolor{computestroke}{Compute}       & Memory-bound kernels    & Operator fusion, FlashAttention \\
  Throughput plateau        & \textcolor{datastroke}{Communication}    & AllReduce saturated     & Gradient compression, overlap \\
  Periodic 30s stalls       & \textcolor{routingstroke}{Coordination}  & Sync checkpointing      & Async checkpointing \\
  Throughput variance       & \textcolor{routingstroke}{Coordination}  & Straggler nodes         & Straggler mitigation \\
  Gradient staleness        & \textcolor{datastroke}{Communication}    & Network congestion      & Traffic engineering, ECMP \\
  Scheduling delays         & \textcolor{routingstroke}{Coordination}  & Resource fragmentation  & Gang scheduling, preemption \\
  \bottomrule
\end{tabular}

\vspace{0.15cm}
\mlsysinsight{Pattern:}{Diagnosis precedes optimization. Always ask ``which C?'' before tuning.}

\end{frame}

% --- ACTIVE LEARNING 1: Predict ---
\begin{frame}{Predict: Where Is the Bottleneck?}
\note{
% -- NARRATE: What to SAY while showing this slide
Read the scenario aloud. ``You are training a 70B model across 4,096
GPUs. Throughput is only 40\% of theoretical peak. Which C3 axis is
most likely the bottleneck?'' Emphasize: 40\% means 60\% of silicon
is doing nothing useful.

% -- ENGAGE: Specific question, prediction, or task for THIS slide
Give 60 seconds for individual writing. Then 30 seconds of pair
discussion. Cold-call 2--3 students. Expected answer: Communication ---
at 4,096 GPUs, AllReduce overhead for 70B parameters (140 GB gradients)
is massive. Some students will say Compute (kernel inefficiency) ---
acknowledge it but redirect: at 40\% peak with 4,096 GPUs, the
Communication term is almost certainly dominant.

% -- FLEX: [CORE] or [OPTIONAL] + contingency
[CORE] This is the first active learning moment --- establishes the
predict-before-reveal pattern for the entire course.
IF SHORT: Reduce to 30 seconds writing, skip pair discussion, cold-call
one student.
}

\centering
\vspace{0.8cm}
{\Large\bfseries Think--Write--Share}

\vspace{0.4cm}
{\large You are training a 70B parameter model across 4,096 GPUs.\\
Training throughput is 40\% of theoretical peak.}

\vspace{0.4cm}
{\normalsize Which \Cthree{} axis is most likely the bottleneck?\\[0.1cm]
Write your answer and one reason why. \textcolor{midgray}{(60 seconds)}}

\vspace{0.3cm}
{\small\textcolor{lightgray}{Compare with a neighbor --- do you agree?}}

\end{frame}

% =============================================================================
\section{The Fleet Stack}
% =============================================================================

\begin{frame}{The Fleet Stack: Four Layers}
\note{
% -- LINK: What prior concept connects to this slide
C3 diagnoses WHERE the bottleneck is. The Fleet Stack organizes HOW the
course addresses each layer of the system --- from physical silicon to
societal governance.

% -- NARRATE: What to SAY while showing this slide
Walk bottom-to-top through the diagram: ``Infrastructure provides the
physical substrate --- compute, network, storage. Distributed ML adds
the training algorithms --- data, tensor, pipeline parallelism.
Deployment puts models into production --- inference optimization,
scheduling, SLAs. Governance ensures responsible operation --- security,
sustainability, fairness.'' Pause on the insight: ``You cannot skip
layers. A governance failure is ultimately an infrastructure failure ---
if the cluster cannot audit which data trained which model, no amount
of policy fixes that.''

% -- ENGAGE: Specific question, prediction, or task for THIS slide
Ask: ``Which layer do most ML courses stop at?'' Expected answer:
Layer 2 (Distributed ML). ``This course covers all four --- production
ML fails when any layer is neglected.''

% -- WARN: What students will get wrong on THIS topic
Students will treat layers as independent. In reality, decisions at the
bottom constrain possibilities at the top: a network topology that
cannot isolate tenants makes multi-tenant governance impossible.
IF STUCK: Give the concrete example: ``If your IB fabric has no SR-IOV,
you cannot do secure multi-tenant serving.''

% -- FLEX: [CORE] or [OPTIONAL] + contingency
[CORE] The Fleet Stack is the course roadmap --- every chapter maps
to a layer.
IF AHEAD: ``Which layer do you think has the highest dollar-cost of
getting wrong?''
IF SHORT: Name the four layers, read the insight callout, move on.
}

% --- Layout: FULL-WIDTH diagram ---
\centering
\safeimg[width=0.92\textwidth,height=4.8cm,keepaspectratio]{fleet-stack.pdf}

\vspace{0.1cm}
\mlsysconcept{Design principle:}{Each layer builds on the one below --- production systems require \textbf{all four}.}

\end{frame}

\begin{frame}{Fleet Stack: What Each Layer Teaches}
\note{
% -- LINK: What prior concept connects to this slide
The previous slide showed the Fleet Stack as a diagram. This table
translates it into concrete skills students will develop in each layer.

% -- NARRATE: What to SAY while showing this slide
Read each row as a promise: ``In the Infrastructure layer, you will
learn to reason about cluster hardware --- GPU selection, InfiniBand
topology, storage hierarchy. In Distributed ML, you will design parallel
training --- DP, TP, PP, collective ops, fault-tolerant training.''
Emphasize the crimson card: ``Most courses stop at Layer 2. This course
covers all four.''

% -- ENGAGE: Specific question, prediction, or task for THIS slide
Ask: ``Which layer interests you most and why?'' Quick show of hands
for each layer. Use the distribution to preview which chapters will
resonate most.

% -- WARN: What students will get wrong on THIS topic
Students will discount Governance as ``not technical.'' Counter with:
the EU AI Act requires auditable training provenance --- that is a
distributed systems problem (data lineage across a 10,000-GPU fleet),
not a policy problem.
IF STUCK: ``If you cannot prove which data trained your model, you
cannot deploy in Europe. That is infrastructure.''

% -- FLEX: [OPTIONAL] + contingency
[OPTIONAL] The Fleet Stack was covered on the previous slide.
IF SHORT: Skip this slide entirely --- the diagram carries the message.
IF AHEAD: Ask which layer they think is most under-invested at real
companies.
}

\scriptsize
\renewcommand{\arraystretch}{1.15}
\begin{tabular}{@{}lp{4cm}p{5.5cm}@{}}
  \toprule
  \textbf{Layer} & \textbf{You Will Learn To...} & \textbf{Key Skills} \\
  \midrule
  \textcolor{computestroke}{\textbf{Infrastructure}} & Reason about cluster hardware & GPU selection, InfiniBand topology, storage hierarchy \\
  \textcolor{datastroke}{\textbf{Distributed ML}} & Design parallel training & DP/TP/PP, collective ops, fault-tolerant training \\
  \textcolor{routingstroke}{\textbf{Deployment}} & Serve models at scale & Inference optimization, fleet scheduling, SLAs \\
  \textcolor{errorstroke}{\textbf{Governance}} & Build responsible systems & Security, sustainability, fairness at fleet scale \\
  \bottomrule
\end{tabular}

\vspace{0.1cm}
\begin{mlsyscard}{crimson}
{\footnotesize Most courses stop at Layer 2. This course covers all four --- production ML fails when \emph{any} layer is neglected.}
\end{mlsyscard}

\end{frame}

% =============================================================================
\section{Course Roadmap}
% =============================================================================

\begin{frame}{17 Chapters in Four Parts}
\note{
% -- LINK: What prior concept connects to this slide
The Fleet Stack named four layers. This roadmap maps those layers to
17 specific chapters across the semester, showing the learning arc.

% -- NARRATE: What to SAY while showing this slide
Point to each part in the diagram: ``Part I is Infrastructure ---
compute, network, storage. Part II is Distributed ML --- parallelism
strategies, collective communication, fault tolerance. Part III is
Deployment --- inference, scheduling, serving. Part IV is Governance ---
security, sustainability, responsible AI.'' Highlight the color coding:
``Notice how C3 threads through all four parts --- Compute constraints
dominate Parts I and II, Communication in Parts II and III, Coordination
in Parts III and IV.''

% -- ENGAGE: Specific question, prediction, or task for THIS slide
Ask: ``Looking at this map, which chapter title are you most curious
about?'' Quick poll. This surfaces student interests early.

% -- WARN: What students will get wrong on THIS topic
Students will assume the parts are independent. In reality, Part III
(Deployment) depends heavily on Part I (Infrastructure) decisions ---
you cannot optimize inference serving without understanding the memory
hierarchy from Chapter 2.
IF STUCK: ``Think of it as a building: you cannot furnish the penthouse
before pouring the foundation.''

% -- FLEX: [OPTIONAL] + contingency
[OPTIONAL] This is a reference slide that students will revisit.
IF SHORT: Show the diagram, name the four parts, move on.
IF AHEAD: Ask students to predict which part will be most relevant
to their career goals.
}

% --- Layout: FULL-WIDTH diagram ---
\centering
\safeimg[width=0.92\textwidth,height=4.8cm,keepaspectratio]{course-roadmap.pdf}

\vspace{0.1cm}
{\small \textcolor{computestroke}{\textbf{Compute}} constraints in Part I \& II \;\textbar\;
\textcolor{datastroke}{\textbf{Communication}} in Part II \& III \;\textbar\;
\textcolor{routingstroke}{\textbf{Coordination}} in Part III \& IV}

\end{frame}

\begin{frame}{The Numbers That Define Fleet Scale}
\note{
% -- LINK: What prior concept connects to this slide
The roadmap showed the course structure. This slide grounds it in
physical reality --- the numbers that make fleet engineering different
from single-node work.

% -- NARRATE: What to SAY while showing this slide
Let the numbers speak. Pause on each one: ``10,000 GPUs. 100 MW of
power --- that is a small city. 350 GB of gradients synchronized every
few seconds. And the one that changes everything: 10,000 GPUs with
100,000h MTBF means one failure every 10 hours. Not one failure per
year. Every 10 hours.''
ANALOGY: ``Imagine a 10,000-person orchestra where one musician
collapses every 10 hours, and the entire orchestra must stop and restart
from a checkpoint.''

% -- ENGAGE: Specific question, prediction, or task for THIS slide
Ask: ``How does this failure rate change how you think about writing
training code?'' Give 20 seconds. Expected insight: checkpointing
becomes the most critical code path, not the model architecture.

% -- WARN: What students will get wrong on THIS topic
Students will try to prevent failures rather than engineer for
resilience. At 10,000 GPUs, prevention is impossible --- the math
guarantees failures. The correct framing is: minimize recovery time,
not failure probability.
IF STUCK: ``You cannot prevent rain. You build a roof.''

% -- FLEX: [CORE] or [OPTIONAL] + contingency
[CORE] These numbers recur throughout the entire course as anchors.
IF AHEAD: ``Calculate the fleet MTBF for 25,000 GPUs with 100,000h
per-GPU MTBF.'' (Answer: 4 hours.)
IF SHORT: Highlight the failure rate number and the bottom-line callout.
}

% --- Layout: FULL-WIDTH diagram ---
\centering
\safeimg[width=0.92\textwidth,height=4.8cm,keepaspectratio]{scale-numbers.pdf}

\vspace{0.1cm}
\mlsysalert{Implication:}{At fleet scale, the ``happy path'' is the exception, not the rule.}

\end{frame}

\begin{frame}{The Warehouse-Scale Computer}
\note{
% -- LINK: Scale Numbers gave abstract numbers; this photo makes them visceral
Students just saw 10,000 GPUs and 100 MW on a slide. This photo shows what
that LOOKS like in the real world --- rows of racks stretching to the horizon.

% -- NARRATE: What to SAY while showing this slide
``This is what 10,000 GPUs looks like. Each rack: 8 GPUs, 40 kW, liquid cooling.
This facility consumes more power than a small town. The network fabric connecting
these racks has more bandwidth than a small country's internet. This is NOT a
server room --- this is a warehouse-scale computer.''

% -- ENGAGE: Specific question for THIS slide
``What is the monthly electricity bill for a 100 MW datacenter at \$0.05/kWh?''
Give 15 seconds. Answer: 100,000 kW $\times$ 720 hours $\times$ \$0.05 = \$3.6M/month.

% -- WARN: Students underestimate physical constraints. Power delivery, cooling,
and floor space are often the binding constraints --- not GPU availability.

% -- FLEX: [CORE] Grounds the abstract numbers in physical reality.
IF SHORT: Show the photo, state the power number, move on.
}

\centering
\safeimg[width=0.88\textwidth,height=4.5cm,keepaspectratio]{datacenter-floor.jpg}

\vspace{0.1cm}
\mlsysalert{Physical reality:}{100 MW of power, liquid cooling per rack, and a network fabric with more bandwidth than a small country's internet.}

\end{frame}

% --- ACTIVE LEARNING: Micro-Retrieval Cue ---
\begin{frame}{Quick Check}
\note{
% -- NARRATE: What to SAY while showing this slide
Read the question aloud: ``If one GPU fails once per year, how often
does a 10,000-GPU cluster fail?'' Pause 15 seconds. Cold-call.
Answer: approximately once per hour (10,000 failures/year divided by
8,760 hours/year is about 1.14 per hour).

% -- FLEX: [CORE] or [OPTIONAL] + contingency
[CORE] This cements the failure-rate intuition before moving on.
IF SHORT: Ask, pause 10 seconds, give the answer, move on.
}

\centering
\vspace{1.0cm}
{\Large\bfseries Quick Check}

\vspace{0.6cm}
{\large If one GPU fails once per year,\\how often does a 10,000-GPU cluster fail?}

\vspace{0.5cm}
{\normalsize\textcolor{midgray}{15 seconds --- then I will cold-call.}}

\end{frame}



% =============================================================================
\section{What You'll Learn}
% =============================================================================

\begin{frame}{Learning Outcomes}
\note{
% -- LINK: What prior concept connects to this slide
The roadmap and scale numbers established what the course covers and
why. This slide translates that into measurable skills students will
demonstrate.

% -- NARRATE: What to SAY while showing this slide
Read through outcomes emphasizing the verbs: ``Design --- not describe,
design. Calculate --- not estimate, calculate. Every outcome is
quantitative. When we say `design distributed training pipelines,' we
mean you will specify TP degree, PP stages, and DP replicas for a given
model and cluster, then calculate the expected MFU.''

% -- ENGAGE: Specific question, prediction, or task for THIS slide
Ask: ``Which of these seven outcomes do you currently feel least
prepared for?'' Quick show of hands per outcome. Use the distribution
to calibrate pacing for early chapters.

% -- WARN: What students will get wrong on THIS topic
Students will underestimate outcome 7 (governance). It requires the
same quantitative rigor as the others --- calculating carbon footprint
per training run, auditing data provenance across a distributed pipeline.
IF STUCK: ``Governance is not an essay. It is a systems design problem
with measurable constraints.''

% -- FLEX: [OPTIONAL] + contingency
[OPTIONAL] Outcomes are a reference --- students will revisit on the
syllabus.
IF SHORT: Skim the list, emphasize the quantitative verbs, move on.
IF AHEAD: Ask students to rank the outcomes by difficulty.
}

\footnotesize
By the end of this course, you will be able to:

\vspace{0.1cm}
\begin{enumerate}\setlength\itemsep{0pt}
  \item \textbf{Design} distributed training pipelines using data, tensor, and pipeline parallelism
  \item \textbf{Calculate} communication costs for collective operations across cluster topologies
  \item \textbf{Reason} about fault tolerance: checkpointing strategies, MTBF analysis, recovery time
  \item \textbf{Optimize} fleet scheduling for utilization, fairness, and SLA compliance
  \item \textbf{Evaluate} inference serving systems for throughput, latency, and cost efficiency
  \item \textbf{Apply} the \Cthree{} framework to diagnose fleet-scale performance bottlenecks
  \item \textbf{Assess} governance requirements: security, sustainability, responsible deployment
\end{enumerate}

\end{frame}

% --- ACTIVE LEARNING: Fermi Estimate Exercise ---
\begin{frame}{Your Turn: How Many GPUs to Train GPT-4?}
\note{
% -- LINK: Learning outcomes promised quantitative reasoning; this exercise
delivers it immediately. Students apply Fermi estimation to fleet-scale ML.

% -- NARRATE: What to SAY while showing this slide
``This is a napkin-math exercise. No calculators. You have the hints on the
right. 90 seconds.'' After pause: ``GPT-4 reportedly used $\sim$2 $\times$ 10$^{25}$
FLOPs. Each H100 delivers $\sim$1,000 TFLOPS (FP16) at $\sim$50\% MFU = 500 TFLOPS
effective. Per day: 500 $\times$ 10$^{12}$ $\times$ 86,400 = 4.3 $\times$ 10$^{19}$
FLOPs/GPU-day. Total GPU-days: 2 $\times$ 10$^{25}$ / 4.3 $\times$ 10$^{19}$
$\approx$ 465,000 GPU-days. For a 90-day run: $\sim$5,000 GPUs. For 30 days:
$\sim$15,000 GPUs. Reported: $\sim$25,000 A100s (MFU was lower).''

% -- ENGAGE: This IS the active learning moment.
After revealing: ``What assumptions dominated your error? Most students
underestimate MFU loss. Real clusters achieve 30--50\% MFU, not 100\%.''

% -- WARN: Students will use peak TFLOPS and get an answer 2--3x too low.
The MFU correction is the key lesson: peak performance is a fiction.

% -- FLEX: [CORE] This exercise sets the quantitative tone for the course.
IF SHORT: Give 60 seconds instead of 90, walk through the solution quickly.
}

\small
\begin{columns}[T]
  \begin{column}{0.55\textwidth}
    {\normalsize\bfseries Fermi Estimate: GPT-4 Training Fleet}

    \vspace{0.15cm}
    GPT-4 required $\sim$$2 \times 10^{25}$ FLOPs of compute.\\
    You have a cluster of H100 GPUs.

    \vspace{0.15cm}
    \textbf{Estimate:}
    \begin{enumerate}\setlength\itemsep{1pt}
      \item How many GPU-days of compute?
      \item How many GPUs for a 90-day training run?
    \end{enumerate}

    \vspace{0.1cm}
    {\footnotesize\textcolor{midgray}{(90 seconds --- napkin math only)}}
  \end{column}
  \begin{column}{0.42\textwidth}
    \textbf{Hints:}
    \begin{itemize}\setlength\itemsep{0pt}
      \item {\footnotesize H100 peak: 1,000 TFLOPS (FP16)}
      \item {\footnotesize Realistic MFU: $\sim$50\%}
      \item {\footnotesize Seconds per day: 86,400}
    \end{itemize}

    \vspace{0.15cm}
    \pause
    \begin{mlsyscard}{computestroke}
    {\scriptsize \textbf{Answer:} $\sim$465K GPU-days.\\
    90-day run $\to$ $\sim$5,000 H100s.\\
    Reported: $\sim$25K A100s (lower MFU).\\[0.05cm]
    \textcolor{computestroke}{\textbf{MFU is the hidden variable!}}}
    \end{mlsyscard}
  \end{column}
\end{columns}

\end{frame}

\begin{frame}{Visual Language}
\note{
% -- NARRATE: What to SAY while showing this slide
Point to each card: ``Blue means compute or processing --- GPU ops,
forward/backward pass. Green means data flow or healthy paths. Orange
means routing or scheduling. Red means error, cost, or bottleneck.
These colors are consistent across every diagram and slide in the
course. When you see red in a figure, something is wrong or expensive.''

% -- FLEX: [OPTIONAL] + contingency
[OPTIONAL] Reference slide for the color system.
IF SHORT: Skip entirely --- students will absorb the colors through
exposure.
}

\small
Throughout this course, colors carry meaning:

\vspace{0.3cm}
\begin{columns}[T]
  \begin{column}{0.45\textwidth}
    \begin{mlsyscard}{computestroke}
    \textbf{Blue} --- Compute / Processing\\
    {\footnotesize GPU ops, forward/backward pass, inference}
    \end{mlsyscard}
    \vspace{0.15cm}
    \begin{mlsyscard}{datastroke}
    \textbf{Green} --- Data / Memory\\
    {\footnotesize Data flow, caches, healthy paths}
    \end{mlsyscard}
  \end{column}
  \begin{column}{0.45\textwidth}
    \begin{mlsyscard}{routingstroke}
    \textbf{Orange} --- Routing / Scheduling\\
    {\footnotesize Load balancers, batch windows}
    \end{mlsyscard}
    \vspace{0.15cm}
    \begin{mlsyscard}{errorstroke}
    \textbf{Red} --- Error / Cost / Bottleneck\\
    {\footnotesize Loss, decode phase, waste}
    \end{mlsyscard}
  \end{column}
\end{columns}
\end{frame}



\begin{frame}{A Taste of What's Coming}
\note{
% -- LINK: The learning outcomes listed abstract skills. This visual roadmap shows
the four-part progression and makes it concrete with a frontier training scenario.

% -- NARRATE: What to SAY while showing this slide
Walk through the roadmap SVG: ``The course follows a four-part arc. Part I builds
the physical infrastructure: compute, network, storage. Part II teaches how to
split training across that infrastructure. Part III deploys trained models at
scale. Part IV addresses governance: security, sustainability, responsibility.''
Then point to the scenario column: ``At each stage, real failure modes arise.
A GPU dies every 4 hours. AllReduce across 25,000 GPUs takes 500 ms per step.
These are not hypothetical --- they are daily realities at frontier labs.''

% -- ENGAGE: Specific question, prediction, or task for THIS slide
Ask: ``Which of these problems would you know how to solve today?''
Expected answer: none of them --- and that gap is exactly why they are
taking this course.

% -- WARN: Students try to solve each problem in isolation. The real challenge
is that these problems interact: a GPU failure triggers checkpointing,
which saturates the network, which causes gradient staleness, which spikes loss.
IF STUCK: ``These are not five separate problems. They are one system
with five failure modes.''

% -- FLEX: [CORE] This is the motivational hook that justifies the entire course.
IF AHEAD: ``What is the dollar cost of 4 hours of 25,000 idle H100s
at \$3/GPU-hour?'' (Answer: \$300K wasted per failure event.)
IF SHORT: Point to the diagram, name the four parts, move on.
}

% --- Layout: FULL-WIDTH visual roadmap ---
\centering
\safeimg[width=0.92\textwidth,height=4.2cm,keepaspectratio]{course-progression.pdf}

\vspace{0.1cm}
\scriptsize
\begin{columns}[T]
  \begin{column}{0.48\textwidth}
    \textbf{What can go wrong at each stage:}
    \begin{itemize}\setlength\itemsep{0pt}
      \item GPU dies every 4h --- restart from scratch?
      \item AllReduce across 25K GPUs: 500ms/step overhead
      \item Rack switch fails --- 128 GPUs go dark
      \item Power budget limits to 80\% utilization
    \end{itemize}
  \end{column}
  \begin{column}{0.48\textwidth}
    \textbf{What this course teaches:}
    \begin{itemize}\setlength\itemsep{0pt}
      \item \textcolor{computestroke}{\textbf{Part I:}} Infrastructure (Ch 2--4)
      \item \textcolor{datastroke}{\textbf{Part II:}} Distributed ML (Ch 5--7)
      \item \textcolor{routingstroke}{\textbf{Part III:}} Deployment (Ch 8--11)
      \item \textcolor{errorstroke}{\textbf{Part IV:}} Governance (Ch 12--17)
    \end{itemize}
  \end{column}
\end{columns}

\vspace{0.05cm}
\centering
{\footnotesize Every failure mode above is a \textbf{chapter in this course}.}

\end{frame}

% --- ACTIVE LEARNING 2: Discussion ---
\begin{frame}{Discussion: What Breaks First?}
\note{
% -- LINK: What prior concept connects to this slide
The scenario slide listed five failure modes. This discussion forces
students to reason about which one strikes first --- building intuition
for failure ordering at scale.

% -- NARRATE: What to SAY while showing this slide
Read the scenario: ``Your startup raised \$100M. You buy 10,000 H100
GPUs, connect them with InfiniBand. What breaks first?'' Point to the
five options along the bottom.

% -- ENGAGE: Specific question, prediction, or task for THIS slide
Turn-and-talk for 90 seconds. Then cold-call 2--3 pairs. There is no
single right answer --- the point is that ``more GPUs'' is never the
full answer. Common student answer: ``the network'' --- press them on
whether that is Communication (bandwidth saturation) or Coordination
(scheduling, straggler management). Some will say ``your budget'' ---
that is a valid and insightful answer worth exploring.

% -- WARN: What students will get wrong on THIS topic
Students will focus on GPU hardware failures. In practice, the software
stack (NCCL hangs, CUDA OOM, driver crashes) breaks far more often than
physical hardware. Meta's Grand Teton paper reports that software issues
cause more downtime than hardware failures.
IF STUCK: ``Think about what you have to set up BEFORE the GPUs start
computing.''

% -- FLEX: [CORE] or [OPTIONAL] + contingency
[CORE] This is the second active learning moment and the most
interactive slide in the deck.
IF SHORT: Reduce to 60-second pairs, cold-call one pair.
IF AHEAD: After discussion, ask: ``What would you monitor to detect
the failure before it happens?''
}

\centering
\vspace{0.8cm}
{\large\bfseries Turn and Talk \textcolor{midgray}{(90 seconds)}}

\vspace{0.5cm}
{\large Your startup just raised \$100M to train a foundation model.\\
You buy 10,000 H100 GPUs and connect them with InfiniBand.\\[0.3cm]
\alert{What breaks first?}}

\vspace{0.5cm}
\begin{columns}[c]
  \begin{column}{0.18\textwidth}\centering\small Network\\ Fabric\end{column}
  \begin{column}{0.18\textwidth}\centering\small GPU\\ Failures\end{column}
  \begin{column}{0.18\textwidth}\centering\small Software\\ Stack\end{column}
  \begin{column}{0.18\textwidth}\centering\small Cooling /\\ Power\end{column}
  \begin{column}{0.18\textwidth}\centering\small Your\\ Budget\end{column}
\end{columns}

\end{frame}

% =============================================================================
\section{Prerequisites \& Context}
% =============================================================================

\begin{frame}{Prerequisites}
\note{
% -- LINK: What prior concept connects to this slide
The course content and motivation are established. Now students need to
know: am I prepared for this?

% -- NARRATE: What to SAY while showing this slide
Walk through the Required column: ``Vol I or equivalent --- you must
understand single-GPU training, the memory hierarchy, and basic
profiling. PyTorch proficiency --- you will write torch.distributed
code from week 2. Linux and SSH --- you will be running jobs on
multi-node clusters.'' Then the Helpful column: ``Distributed systems
concepts like consensus and RPC will be introduced as needed. If you
know Slurm or Kubernetes, you will have a head start on the scheduling
chapters, but it is not required.''

% -- ENGAGE: Specific question, prediction, or task for THIS slide
Ask: ``Raise your hand if you are comfortable explaining what AllReduce
does.'' The fraction of hands up calibrates how much Ch.~5--6 review
is needed.

% -- WARN: What students will get wrong on THIS topic
Students without cluster experience will feel behind. Reassure them:
the course introduces cluster concepts from scratch. The real
prerequisite is single-GPU fluency, not distributed systems expertise.
IF STUCK: ``If you can train a ResNet on one GPU and profile where time
goes, you are ready.''

% -- FLEX: [OPTIONAL] + contingency
[OPTIONAL] Logistical slide.
IF SHORT: Name the top 3 required skills, mention the ``helpful but not
required'' list exists, move on.
IF AHEAD: Ask a student with cluster experience to share one surprising
lesson from their first multi-node job.
}

\small
\begin{columns}[T]
  \begin{column}{0.48\textwidth}
    \textbf{Required:}
    \begin{itemize}\setlength\itemsep{1pt}
      \item \textbf{Vol I} or equivalent ML systems foundation
      \item Single-GPU training and inference experience
      \item PyTorch proficiency (\texttt{torch.nn}, \texttt{DataLoader})
      \item Basic linear algebra and probability
      \item Linux command line and SSH
    \end{itemize}
  \end{column}
  \begin{column}{0.48\textwidth}
    \textbf{Helpful but not required:}
    \begin{itemize}\setlength\itemsep{1pt}
      \item Distributed systems concepts (consensus, RPC)
      \item Networking basics (TCP/IP, bandwidth, latency)
      \item Cluster job schedulers (Slurm, Kubernetes)
      \item CUDA programming
    \end{itemize}

    \vspace{0.15cm}
    {\footnotesize\textcolor{midgray}{We will introduce distributed systems concepts as needed throughout the course.}}
  \end{column}
\end{columns}

\end{frame}

\begin{frame}{Relationship to Volume I}
\note{
% -- LINK: What prior concept connects to this slide
Prerequisites established what students need. This slide addresses the
elephant in the room: ``Is this harder than Vol I?''

% -- NARRATE: What to SAY while showing this slide
Point to the two cards side by side: ``Vol I: one machine, 1--8 GPUs,
shared memory, DataParallel, the Iron Law. Vol II: 10,000+ GPUs,
InfiniBand, message passing, torch.distributed, C3.'' Then the key
message: ``The distinction is SCOPE, not DEPTH. Vol I went deep on one
machine. Vol II goes wide across the fleet. Both are equally rigorous.
We re-derive key frameworks like the Iron Law at fleet scale, adding
communication and coordination terms.''

% -- ENGAGE: Specific question, prediction, or task for THIS slide
Ask: ``What Vol I concept do you think will change most at fleet scale?''
Cold-call. Any answer works --- use it to preview how that concept
evolves. (e.g., ``the memory wall'' becomes the ``network wall.'')

% -- WARN: What students will get wrong on THIS topic
Students will assume Vol II content is ``harder.'' It is not --- it is
different. A student who struggled with cache hierarchies in Vol I may
excel at distributed fault tolerance in Vol II. The thinking patterns
are complementary, not hierarchical.
IF STUCK: ``Think of it as learning to fly after learning to drive.
Different skills, not harder skills.''

% -- FLEX: [OPTIONAL] + contingency
[OPTIONAL] Context-setting slide.
IF SHORT: Read the crimson card at the bottom, skip the detailed
comparison.
IF AHEAD: Ask: ``Which Vol I equation do you think we will generalize
first?'' (Answer: the Iron Law, in Chapter 1.)
}

\small
\begin{columns}[T]
  \begin{column}{0.48\textwidth}
    \begin{mlsyscard}{computestroke}
    \textbf{Vol I: The Machine}\\[0.1cm]
    {\footnotesize
    \begin{itemize}\setlength\itemsep{0pt}
      \item Single node, 1--8 GPUs
      \item PCIe / NVLink within one box
      \item Shared memory model
      \item \texttt{torch.nn.DataParallel}
      \item Iron Law: $T = D/BW + O/R + L$
    \end{itemize}
    }
    \end{mlsyscard}
  \end{column}
  \begin{column}{0.48\textwidth}
    \begin{mlsyscard}{datastroke}
    \textbf{Vol II: The Fleet}\\[0.1cm]
    {\footnotesize
    \begin{itemize}\setlength\itemsep{0pt}
      \item Multi-node, 10,000+ GPUs
      \item InfiniBand / Ethernet fabric
      \item Message passing model
      \item \texttt{torch.distributed}
      \item \Cthree{}: Compute + Comm + Coord
    \end{itemize}
    }
    \end{mlsyscard}
  \end{column}
\end{columns}

\vspace{0.2cm}
\centering
\begin{mlsyscard}{crimson}
{\footnotesize\textbf{Scope, not depth.} Both volumes are equally rigorous. Vol II re-establishes key frameworks --- like the Iron Law --- at fleet scale, adding communication and coordination terms.}
\end{mlsyscard}

\end{frame}

\begin{frame}{What Changes at 1,000 GPUs?}
\note{
% -- LINK: Vol I comparison showed the two scopes; this slide bridges them
by naming the specific qualitative shifts that occur at fleet scale.

% -- NARRATE: What to SAY while showing this slide
``Vol I taught you to think about one machine. Here are five things that
change qualitatively --- not just quantitatively --- when you cross the
node boundary. The bus becomes a network. Shared memory becomes message
passing. Rare failures become routine. Profiling becomes distributed
tracing. And cost shifts from hardware to operations.''

% -- ENGAGE: Specific question for THIS slide
``Which of these five shifts surprises you most?'' Cold-call one student.
Use their answer to preview how the course addresses it.

% -- WARN: Students think fleet-scale is just ``more of the same.'' Correct:
each shift introduces qualitatively new failure modes and design patterns.

% -- FLEX: [CORE] This is the bridge from Vol I to Vol II.
IF SHORT: Name all five, spend time on the failure rate shift only.
}

\small
\centering
\renewcommand{\arraystretch}{1.3}
{\footnotesize
\begin{tabular}{@{}lll@{}}
  \toprule
  \textbf{Dimension} & \textbf{1--8 GPUs (Vol I)} & \textbf{1,000+ GPUs (Vol II)} \\
  \midrule
  \textbf{Interconnect} & NVLink: 900 GB/s & InfiniBand: 50 GB/s (18$\times$ slower) \\
  \textbf{Memory model} & Shared memory & Message passing (explicit sync) \\
  \textbf{Failure rate} & MTBF $\sim$years & MTBF $\sim$hours (routine) \\
  \textbf{Debugging} & \texttt{nvidia-smi}, profiler & Distributed tracing, fleet telemetry \\
  \textbf{Cost driver} & GPU hardware (\$) & Operations + idle time (\$\$\$) \\
  \bottomrule
\end{tabular}
}

\vspace{0.2cm}
\mlsysalert{Key insight:}{Fleet-scale is not ``more of the same.'' Each row is a \textbf{qualitative} shift that demands new abstractions.}

\end{frame}

% =============================================================================
\section{Wrap-Up}
% =============================================================================

\begin{frame}{Fallacies}
\note{
% -- NARRATE: Three common misconceptions about fleet-scale ML courses.
Read each fallacy, then deliver the quantitative rebuttal.
Spend extra time on the first: students consistently overestimate
how much single-node experience transfers to fleet engineering.

% -- FLEX: [CORE] Fallacies set expectations before the course begins.
IF SHORT: Cover the first two fallacies, skip the third.
}

\small
\textbf{Fallacy:} \textit{Fleet-scale ML is just single-node ML with more GPUs.}\\
{\footnotesize At 1,000 GPUs, communication consumes 30--40\% of iteration time. Failure recovery, not model architecture, dominates engineering effort. Qualitatively different.}

\vspace{0.15cm}
\textbf{Fallacy:} \textit{You need distributed systems experience before taking this course.}\\
{\footnotesize Vol I foundations (memory hierarchy, profiling, single-GPU training) are the real prerequisites. Distributed concepts (consensus, AllReduce) are taught from scratch.}

\vspace{0.15cm}
\textbf{Fallacy:} \textit{Governance chapters are ``soft'' --- the real course is infrastructure.}\\
{\footnotesize Calculating carbon footprint per training run, auditing data provenance, and designing privacy-preserving pipelines require the same quantitative rigor as sizing a network fabric.}

\end{frame}

\begin{frame}{Pitfalls}
\note{
% -- NARRATE: Three operational mistakes students make early in the course.
% -- FLEX: [CORE] IF SHORT: Cover the first pitfall only.
}

\small
\textbf{Pitfall:} \textit{Skipping the napkin math.}\\
{\footnotesize Every design decision in this course starts with a back-of-envelope calculation. A 2-minute Fermi estimate prevents weeks of wasted GPU time. Build the habit from day one.}

\vspace{0.15cm}
\textbf{Pitfall:} \textit{Treating peak TFLOPS as achievable throughput.}\\
{\footnotesize Real clusters achieve 30--50\% of peak (MFU). Designing for peak means under-provisioning by 2--3$\times$.}

\vspace{0.15cm}
\textbf{Pitfall:} \textit{Optimizing components in isolation.}\\
{\footnotesize A faster GPU without a faster network increases idle time. A faster network without better scheduling wastes bandwidth. Fleet engineering is co-design.}

\end{frame}

% --- ACTIVE LEARNING 3: Retrieval Practice ---

% --- MUDDIEST POINT ---
\begin{frame}{Muddiest Point}
\note{
% -- NARRATE: What to SAY while showing this slide
``Before we close, I want to know what confused you most. Write one
sentence --- the concept you found muddiest. Anonymous. Submit before
you leave --- slip of paper or digital form.'' Scan responses after
class and address the top 2--3 confusions at the start of next lecture.

% -- FLEX: [CORE] or [OPTIONAL] + contingency
[CORE] Closes the feedback loop --- essential for adaptive teaching.
IF SHORT: Reduce to ``write one word on a slip of paper.''
}

\centering
\vspace{1.0cm}
{\Large\bfseries What was the \alert{muddiest point} today?}

\vspace{0.8cm}
{\normalsize Write down the concept you found \textbf{most confusing}.}

\vspace{0.5cm}
{\small\textcolor{midgray}{Anonymous. One sentence. Submit before you leave.}}

\end{frame}

\begin{frame}{What Were the Key Ideas?}
\note{
% -- NARRATE: What to SAY while showing this slide
``Close your notes. 90 seconds. Write down the 3 most important ideas
from today. No peeking.'' Walk around the room while students write.
Do NOT show the next slide yet --- the struggle to recall is where
learning happens.

% -- FLEX: [CORE] or [OPTIONAL] + contingency
[CORE] Retrieval practice is the highest-impact learning technique
(Rosenshine). Never skip.
IF SHORT: Reduce to 60 seconds.
}

\centering
\vspace{1.5cm}
{\Large\bfseries Close your notes.}

\vspace{0.8cm}
{\large Write down the \textbf{3 most important ideas} from today.}

\vspace{0.8cm}
{\normalsize\textcolor{midgray}{90 seconds --- no peeking.}}

\end{frame}

\begin{frame}{Key Takeaways}
\note{
% -- LINK: What prior concept connects to this slide
Students just attempted recall. This slide reveals the answers and
fills gaps.

% -- NARRATE: What to SAY while showing this slide
Walk through each bullet, pausing on the quantitative anchors: ``The
fleet is the computer --- thousands of accelerators as one unit. Scale
is qualitative --- not just more, but fundamentally different. C3
framework --- every bottleneck maps to Compute, Communication, or
Coordination. Fleet Stack --- four layers from infrastructure to
governance. Failure math --- one failure every 10 hours at 10,000 GPUs.
Scope not depth --- Vol II is wider, not harder.''

% -- FLEX: [CORE] or [OPTIONAL] + contingency
[CORE] Consolidation slide.
IF SHORT: Read just the first three bullets.
}

\scriptsize
\begin{itemize}\setlength\itemsep{0pt}
  \item \textbf{The fleet is the computer:} The unit of design is thousands of coordinated accelerators.
  \item \textbf{Scale is qualitative:} Communication dominates compute, failure becomes routine, behavior emerges.
  \item \textbf{\Cthree{} framework:} Every bottleneck maps to Compute, Communication, or Coordination.
  \item \textbf{Fleet Stack:} Four layers --- Infrastructure, Distributed ML, Deployment, Governance.
  \item \textbf{Failure math:} 10,000 GPUs $\times$ 100,000h MTBF = one failure every 10 hours.
  \item \textbf{Scope, not depth:} Vol II is not ``harder'' than Vol I --- it is \emph{wider}. Equally rigorous.
  \item \textbf{17 chapters, 4 parts:} Infrastructure to governance --- the full production ML lifecycle.
\end{itemize}

\end{frame}

\begin{frame}{References}
\note{
% -- NARRATE: What to SAY while showing this slide
``These five references are the foundational readings for the course.
Verbraeken is the survey that maps the landscape. Narayanan and Jiang
show how frontier labs actually train at scale. Dean is the historical
origin. Patterson and Hennessy is the pedagogical model we follow.''

% -- FLEX: [OPTIONAL] + contingency
[OPTIONAL] Reference slide.
IF SHORT: Skip verbal walkthrough --- students can read it.
}

\small
\mlsysref{Verbraeken+20}{Verbraeken et al. ``A Survey on Distributed ML.'' ACM Computing Surveys, 2020.}
\mlsysref{Narayanan+21}{Narayanan et al. ``Efficient Large-Scale Language Model Training on GPU Clusters.'' SC 2021.}
\mlsysref{Jiang+24}{Jiang et al. ``MegaScale: Scaling Model Training to 10,000+ GPUs.'' NSDI 2024.}
\mlsysref{Dean12}{J. Dean. ``Large-Scale Distributed Systems for Training Neural Networks.'' NIPS 2012.}
\mlsysref{P\&H21}{Patterson \& Hennessy. \textit{Comp.\ Arch.: A Quantitative Approach}. 6th ed.}

\end{frame}

\begin{frame}{Next Lecture: The Distributed Landscape}
\note{
% -- LINK: What prior concept connects to this slide
Today established why scale matters and introduced C3 and the Fleet
Stack. The next lecture dives into the first concrete question: what
is actually inside a modern GPU cluster?

% -- NARRATE: What to SAY while showing this slide
``The fleet is the computer. But what hardware makes up this computer?
Next lecture: Chapter 1 introduces the fleet as a system --- what
accelerators, what interconnects, what topology, and why the answer to
`how do 10,000 GPUs talk to each other?' is the central puzzle of
fleet engineering.'' Point to the three columns: Compute, Communication,
Coordination --- the C3 lens applied to hardware.

% -- FLEX: [CORE] or [OPTIONAL] + contingency
[CORE] Forward hook that creates anticipation.
IF SHORT: Read the central question and move on.
}

\small
\centering
\vspace{0.5cm}
{\Large\bfseries Chapter 1: Introduction to Fleet-Scale ML}

\vspace{0.5cm}
{\large How do 10,000 GPUs coordinate to train one model?}

\vspace{0.3cm}
\begin{columns}[c]
  \begin{column}{0.30\textwidth}
    \centering
    \textcolor{computestroke}{\textbf{Compute}}\\[0.1cm]
    {\footnotesize What hardware?\\Which accelerators?}
  \end{column}
  \begin{column}{0.30\textwidth}
    \centering
    \textcolor{datastroke}{\textbf{Communication}}\\[0.1cm]
    {\footnotesize What topology?\\How fast?}
  \end{column}
  \begin{column}{0.30\textwidth}
    \centering
    \textcolor{routingstroke}{\textbf{Coordination}}\\[0.1cm]
    {\footnotesize What breaks?\\How to recover?}
  \end{column}
\end{columns}

\vspace{0.4cm}
{\normalsize\textbf{The fleet is the computer. Let's understand the hardware.}}

\end{frame}



\appendix

\begin{frame}{Backup: Fleet Scale Quick Reference}
\note{
% -- NARRATE: Backup slide with the key numbers students need for problem sets.
Show if students ask for a reference card during exercises.
% -- FLEX: [OPTIONAL] Only show if requested.
}

\footnotesize
\textbf{Key Fleet-Scale Numbers (reference for problem sets):}

\vspace{0.15cm}
\renewcommand{\arraystretch}{1.15}
\begin{tabular}{@{}lr@{}}
  \toprule
  \textbf{Metric} & \textbf{Value} \\
  \midrule
  H100 peak TFLOPS (FP16) & 1,979 \\
  H100 HBM bandwidth & 3,350 GB/s \\
  NVLink 4.0 (intra-node) & 900 GB/s \\
  InfiniBand NDR (inter-node) & 400 Gbps = 50 GB/s \\
  Typical fleet MFU & 30--50\% \\
  Per-GPU MTBF (typical) & 50,000--100,000 h \\
  10K-GPU fleet MTBF & 5--10 h \\
  GPT-3 gradient size (FP16) & $\sim$700 GB \\
  GPT-4 training compute & $\sim$2 $\times$ 10$^{25}$ FLOPs \\
  \bottomrule
\end{tabular}

\vspace{0.1cm}
{\footnotesize See textbook Appendix A for the complete hardware constants table.}

\end{frame}

\begin{frame}{Backup: Prerequisite Self-Assessment}
\note{
% -- NARRATE: Use if students want to check their readiness before next lecture.
% -- FLEX: [OPTIONAL] Hand out as a take-home diagnostic.
}

\footnotesize
\textbf{Can you answer these Vol I review questions?}

\vspace{0.15cm}
\begin{enumerate}\setlength\itemsep{2pt}
  \item What is the arithmetic intensity of a matrix multiply $C = AB$ where $A$ is $M \times K$ and $B$ is $K \times N$?
  \item A model has 7B parameters in FP16. How much memory for inference? For training with Adam?
  \item What does \texttt{torch.nn.DataParallel} replicate across GPUs: the model, the data, or both?
  \item Your GPU profile shows 80\% of time in memory transfers. Is this workload compute-bound or memory-bound?
  \item What is Amdahl's Law? If 10\% of your code is serial, what is the maximum speedup with 1,000 GPUs?
\end{enumerate}

\vspace{0.1cm}
{\footnotesize If you cannot answer 4 of 5, review Vol I Chapters 1--4 before next lecture.}

\end{frame}


\end{document}
